% !PRINTED 347
\weberpage{347}

%<<CONT
und bestimmen das Vorzeichen so, daß $\pm q \equiv 1 \ (\mathrm{mod}\ 4)$ wird.
Es ist in diesem Falle
\[
\theta = \frac{-1 + \sqrt{\pm q}}{2},
\]
und diesen Ausdruck können wir durch $q$te Einheitswurzeln darstellen. Es bedeute $r$ eine imaginäre $q$te Einheitswurzel und es
durchlaufe $a$ die quadratischen Reste, $b$ die Nichtreste von $q$.
Wir setzen, wie in Bd.\ I, \S\,179
\weq{4}{A = \overset{a}{\Sigma}\, r^a, \quad B = \overset{b}{\Sigma}\, r^b}
und haben nach den dortigen Formeln (3), die ohne Benutzung
des Reziprozitätsgesetzes abgeleitet sind:
\weq{5}{\theta = A, \quad \theta' = B.}
Das Vorzeichen von $\sqrt{+q}$ hängt von der Wahl von $r$ ab,
kommt aber hier nicht in Betracht. Nun ist nach dem polynomischen
 Lehrsatz:
\[
A^p \equiv \Sigma r^{pa}, \quad B^p \equiv \Sigma r^{pb},
\]
also
\[
A^p \equiv A, \quad B^p \equiv B, \quad \text{wenn} \ \left(\frac{p}{q}\right) = +1,
\]
und
\[
A^p \equiv B, \quad B^p \equiv A, \quad \text{wenn} \ \left(\frac{p}{q}\right) = -1.
\]

Da $A$ nicht kongruent $B$ ist, wenn $p$ von $q$ verschieden ist,
so folgt hieraus und aus 2. des vorigen Paragraphen:

3) \gesperrt{Die Kongruenz}
\[
x^2 \equiv \pm q \ (\mathrm{mod}\ p)
\]
\gesperrt{hat dann und nur dann Lösungen, wenn $p$ quadratischer
Rest von $q$ ist.}

Und dies ist das Reziprozitätsgesetz mit seinen Ergänzungssätzen.

\S\,95. \gesperrt{Äquivalente Formen und Ideale im quadratischen Körper.}

Wenn die beiden Ideale [\S\,91, (16)]
\weq{1}{\mathfrak{a} = (a,\, a\,\omega), \quad \mathfrak{a}' = (a',\, a'\,\omega')}
äquivalent sind, so gibt es eine (ganze oder nicht ganze) Zahl $\eta$
in $\Omega$, für die $\eta\, \mathfrak{a}' = \mathfrak{a}$ ist, und diese Zahl $\eta$ ist durch die Ideale
$\mathfrak{a}$, $\mathfrak{a}'$ selbst bis auf einen Faktor, der eine \gesperrt{numerische Einheit}
%CONT>>
